\begin{abstract}
Detailed balance removes directional cycle flow from a Markov model. We evaluate the restriction for next-session open-to-close return terciles of SPY, QQQ and IWM, with 2,766 forecasts per asset over 2015--2025 and a fixed intraday allocation. Nested benchmarks split the unrestricted model's log-loss gain over IID exactly into middle-versus-outer persistence, captured by a two-state chain, additional middle-tercile information and conditional direction. Only QQQ's gain survives every familywise adjustment. This persistence accounts for 82\%, 82\% and 66\% of the gain. Two properties limit this layer: the reversible projection preserves aggregate stationary flow between the middle state and the pooled outer states, and re-estimated terciles balance the marginal direction, keeping the two-state benchmark allocation within about $2\times10^{-4}$ of half exposure. Essentially all incremental net return over IID, 0.043, 0.020 and 0.073 basis points per day, arises in the direction layer, whose pointwise log-loss interval excludes zero only for IWM under the primary rule. Across four estimation rules, its pointwise net-return interval excludes zero in nine of twelve cases, its log-loss interval in one. Reversible-minus-unrestricted log-loss differences (0.0013, 0.0007 and 0.0016 nats; intervals containing zero) are of the size of a closed-form projection cost, which, after a first-order estimation adjustment, agrees with the simulated mean differences within 1.3 Monte Carlo standard errors. Daily liquidation of a near-half exposure costs about five basis points per session, a baseline common to all nine allocations, which all lost capital over the evaluation period; incremental cost differences are below 0.005 basis points per day. The gain is concentrated in 2015--2018 and reverses for SPY and QQQ in 2023--2025. The layers were specified after the primary results.
\end{abstract}

\noindent\textbf{Keywords:} Markov chains; detailed balance; probabilistic forecasting; forecast decomposition; information projection; intraday returns; transaction costs.

\section{Introduction and related research}
Can a structural restriction improve financial probability forecasts, and does any improvement matter for an investment decision? Detailed balance offers a tractable setting for these questions. In a three-state Markov chain it removes one directional circulation parameter while retaining state dependence. The restriction can reduce estimation variance, but it can also discard predictive information when the data-generating process is irreversible. This study evaluates that trade-off using next-session open-to-close return terciles and a fixed allocation rule for three US equity ETFs.

Statistical forecast accuracy and economic value answer different questions \citep{Granger2000,Diebold1995}. Proper scores evaluate the full probability vector \citep{Gneiting2007}, whereas an allocation uses a particular function of that vector and is also exposed to return magnitudes and implementation costs. Predictability and economic value can vary over time \citep{Pesaran1995}; small statistical improvements can be useful under some portfolio mappings \citep{Campbell2008}. A negative trading result under one mapping therefore does not invalidate probabilistic forecasting generally.

The holding interval makes this distinction especially relevant. Trading and non-trading ETF returns differ \citep{Kelly2011}; opening-price pressure, overnight--intraday return patterns and time-of-day variation in risk premia are documented in related markets \citep{Berkman2012,Lou2019,Hendershott2020,Bondarenko2023}. These studies motivate precise entry and exit timing. They do not imply a universal negative intraday premium or determine the execution cost of the ETFs considered here.

Time irreversibility also has an established econometric literature, including time-domain, frequency-domain and distribution-based tests \citep{Ramsey1996,Hinich1998,Chen2000}. Financial applications connect irreversibility with volatility asymmetry and alternative time-series representations \citep{Chen2002,Zumbach2009,Lacasa2016}. Finite-state reversible/irreversible decompositions and reversible estimation methods are already available \citep{McCausland2007,Trendelkamp2015}, and information projections onto reversible kernels have closed forms \citep{WolferWatanabe2021}. The contribution here is an empirical evaluation of a restriction, not a new estimator, projection result or test of market-wide reversibility.

The transferable contribution is an exact accounting that ties score components to the decision rule. The log-loss chain rule separates middle-versus-outer prediction from conditional upper-versus-lower prediction. Two nested benchmarks that share the IID model's marginal direction turn this separation into day-by-day layers, so that the unrestricted-minus-IID difference splits exactly into middle-versus-outer persistence, additional middle-tercile information and conditional direction, for log loss and for net return alike. Two elementary properties then determine what each layer can do. The reversible projection preserves every flow between the middle state and the outer pair, so detailed balance leaves the lumped middle-versus-outer process unchanged at the population level; and re-estimated historical terciles balance upper and lower labels, so a forecast without state-dependent direction holds the allocation at half exposure. The tools are established; the application shows that, in these data, most of the score gain lies where neither the restriction nor the allocation can act. It is not a general decomposition theorem for every proper score and decision rule, and it does not identify a causal mechanism for the economic outcome.

The design links four comparisons. First, identical historical information within each asset isolates the forecast difference between IID, unrestricted Markov and reversible Markov models. Second, the benchmark layers show which probability components improve and whether the allocation can respond to them. Third, a daily-return identity separates fixed half exposure, the probability-driven adjustment and trading cost. Fourth, a closed-form projection cost provides a first-order reference for the size of the restriction's effect, which is compared with the observed contrasts and with a controlled simulation. Circulation paths, alternative estimation rules and calendar segments help interpret the results. SPY is the primary application; QQQ and IWM are related-asset extensions rather than independent replications.

The unrestricted chain ranks first on both scores for all three assets, but the statistical evidence depends on the asset and adjustment method, and most of the gain is middle-versus-outer persistence that a two-state chain reproduces. The incremental economic gain comes from the conditional direction layer, whose score contribution is detectable only for IWM. Detailed balance provides no demonstrated log-loss advantage; its observed cost is of the size the projection reference implies, which is about one standard error of the contrast. Under the fixed daily round-trip rule, small incremental gains coexist with large absolute losses at the primary cost. The full-period score ranking also masks deterioration in the later calendar segment. The benchmark layers, reconstructed-panel checks and simulation were added after the primary results and provide exploratory interpretation of the completed experiment.

\section{Design and estimators}\label{sec:design}
\subsection{Data, timing and state construction}
Adjusted daily Open and Close observations were retrieved on 13 November 2026 from Yahoo Finance through yfinance 1.8.0. Table~\ref{tab:data} describes available histories. The final sample is 2 January 2015--31 December 2025 for all assets. Session coverage matches an XNYS reference calendar; this checks dates, not price accuracy or executable fills.

\begin{equation}
r_t=\frac{C_t}{O_t}-1,\qquad
(q_{1,t},q_{2,t})=\operatorname{Quantile}_{1/3,\,2/3}(r_1,\ldots,r_t).
\label{eq:returns}
\end{equation}
After close $t$, historical returns are relabeled using these cutpoints, with lower state $r\le q_{1,t}$, middle state $q_{1,t}<r\le q_{2,t}$, and upper state $r>q_{2,t}$. They are relative terciles, not absolute losses, flat returns and gains. The target $r_{t+1}$ uses the same saved cutpoints. Quantiles use NumPy's default linear interpolation. Equality at either cutpoint is assigned to the lower of the adjacent states, so ties can prevent exact equal-frequency bins. Write $s_{u\mid t}\in\{1,2,3\}$ for the label of historical return $r_u$ using the cutpoints available at origin $t$. A historical return can change label at a later origin. The experiment therefore evaluates an updating discretization and forecasting procedure, not a single time-homogeneous chain fitted to permanently fixed labels.

An expanding history requires 2,520 observations. SPY development ends in 2009 and validation covers 2010--2014. Protocol v1.1 was internally frozen before final scoring; this was not external preregistration. Past final outcomes enter later fits only after realization. QQQ/IWM apply the same specification sequentially after SPY. No settings were changed in response to final results. This internal freeze limits the reported design choices but is not an independently witnessed safeguard against all selection effects \citep{White2000,Bailey2014}.

At the first final-test origin, SPY, QQQ and IWM have 5,522, 3,980 and 3,672 historical daily observations, respectively: all exceed the minimum, but their initial expanding windows are not equal in length.

\subsection{Forecast estimators}
Let $N_{ij,t}=\sum_{u=1}^{t-1}\mathbf{1}\{s_{u\mid t}=i,s_{u+1\mid t}=j\}$ count adjacent historical transitions and let $\alpha=1$. The IID forecast uses destination labels $s_{2\mid t},\ldots,s_{t\mid t}$:
\begin{equation}
p^{\mathrm{IID}}_{t,j}=\frac{\sum_{i=1}^3N_{ij,t}+\alpha}{t-1+3\alpha}.
\label{eq:iid}
\end{equation}
It is nearly uniform because of the training terciles. The unrestricted estimator is
\begin{equation}
\widehat P_{ij,t}=\frac{N_{ij,t}+\alpha}{\sum_{k=1}^3(N_{ik,t}+\alpha)}.
\label{eq:markov}
\end{equation}
The Markov forecast selects row $s_{t\mid t}$ of the estimated transition matrix. Reversible estimation maximizes
\begin{equation}
\begin{gathered}
\sum_{i,j}(N_{ij,t}+\alpha)\log P_{ij,t},\\
\text{subject to }X=X^\top,\quad X_{ij}>0\ \text{for all }i,j,\quad \sum_{i,j}X_{ij}=1,\\
\pi_i=\sum_jX_{ij},\quad P_{ij}=X_{ij}/\pi_i.
\end{gathered}
\label{eq:reversible}
\end{equation}
This enforces $\pi_iP_{ij}=\pi_jP_{ji}$. Six softmax masses parameterize the three diagonal flows and three off-diagonal pairs; each pair mass is split equally between its two entries. Fixing one logit to zero leaves five free parameters. BFGS uses an analytic gradient, a maximum of 500 iterations and gradient tolerance $10^{-8}$. Initialization uses smoothed diagonal counts and summed smoothed off-diagonal pair counts. A fit is accepted when the optimizer reports success or the final gradient infinity norm is at most $10^{-6}$; non-finite or unaccepted fits raise an error. These criteria document the numerical implementation without proving global optimality. Appendix B distinguishes constraint residuals, optimization diagnostics and independent refitting.

The positive unrestricted family has six free parameters, versus five under reversibility. With $J_{ij}=\pi_iP_{ij}-\pi_jP_{ji}$, stationarity leaves one independent three-state circulation:
\begin{equation}
J_{12}=J_{23}=J_{31};\qquad\text{detailed balance imposes }J_{12}=0.
\label{eq:flow}
\end{equation}
Here $\pi$ is the stationary distribution of the fitted matrix at a given origin; it need not describe a stable distribution throughout the market sample. This restriction removes directional cycle flow, not all conditional asymmetry: $P_{ij}$ and $P_{ji}$ can still differ when stationary masses differ. Crash/rebound narratives are possible motivations, but terciles discard crash magnitude and this study does not identify that mechanism.

\section{Scores and implementation}
\subsection{Probability scores and the information used by the allocation}
For target label $y_{t+1}$ and forecast $p_t$, losses are
\begin{equation}
\ell_{t+1}=-\log p_{t,y_{t+1}},\qquad
b_{t+1}=\sum_{j=1}^3\left(p_{t,j}-\mathbf{1}\{y_{t+1}=j\}\right)^2.
\label{eq:scores}
\end{equation}
Calibration diagnostics use ten fixed bins per class. Empty bins supply no evidence, and bin averages do not display the full range of forecasts. Appendix~\ref{app:families} reports a count-weighted error, a block-bootstrap interval and a simulation reference that assumes independent outcomes conditional on the frozen probability sequence. That reference is more restrictive than calibration alone.

To distinguish the information rewarded by the score from that used in the allocation, define $m_t=p_{t,2}$ and $q_t=p_{t,3}/(p_{t,1}+p_{t,3})$. Let $z_{t+1}=\mathbf{1}\{y_{t+1}=2\}$ and $u_{t+1}=\mathbf{1}\{y_{t+1}=3\}$. Positivity from smoothing makes these quantities well defined. The standard probability chain rule gives an exact sequential binary factorization of categorical log loss \citep[Chapter~2]{CoverThomas1991}:
\begin{equation}
\begin{split}
\ell_{t+1}={}&-z_{t+1}\log m_t-(1-z_{t+1})\log(1-m_t)\\
&-(1-z_{t+1})\{u_{t+1}\log q_t+(1-u_{t+1})\log(1-q_t)\}.
\end{split}
\label{eq:scorefactor}
\end{equation}
The first line evaluates middle versus outer terciles; the second evaluates upper versus lower conditional on an outer-tercile outcome. Meanwhile,
\begin{equation}
p_{t,3}-p_{t,1}=(1-m_t)(2q_t-1).
\label{eq:tiltfactor}
\end{equation}
Improving the middle-versus-outer forecast need not improve upper-versus-lower discrimination. In particular, if $q_t=1/2$, changing $m_t$ does not change the allocation. These identities identify possible channels; Section~\ref{sec:components} measures their empirical contributions from the saved final forecasts. Outer-tercile membership is relative to historical cutpoints, so neither component is itself a test of return sign or a direct measure of volatility.

\subsection{Allocation and transaction costs}
The fixed exposure, chosen after close $t$, is
\begin{equation}
w_t=\operatorname{clip}\left[0.5+0.5(p_{t,3}-p_{t,1}),0,1\right].
\label{eq:weight}
\end{equation}
For valid probability vectors the expression inside the clip already lies in $[0,1]$; clipping is retained to describe the frozen implementation. The position enters at open $t+1$, exits at close $t+1$, and earns zero cash interest:
\begin{equation}
G_{t+1}=w_tr_{t+1},\qquad T_{t+1}=2w_t,\qquad
R_{t+1}(c)=G_{t+1}-cT_{t+1}.
\label{eq:net}
\end{equation}
The primary cost is $c=0.0005$ (5 basis points, bps) per side, with sensitivities at 1, 10 and 25 bps. Costs use starting notional, omitting exact exit-notional adjustment, impact and auction slippage. Approximately half exposure implies approximately one unit of turnover daily, even if consecutive weights are unchanged. Even an uninformative forecast therefore incurs the cost of maintaining the baseline position. Within the observed exposure range, total cost is dominated by that baseline; this need not hold for an arbitrary linear mapping or a different exposure range. The experiment compares models under this rule and does not test whether reversibility has economic value under a thresholded or no-trade policy. Such a policy would need to change participation or explicitly allow positions to persist; an overnight-held alternative would also change the return target and risk exposure.

Define $\delta_t=w_t-1/2$. Under the same accounting, the net return separates exactly into
\begin{equation}
R_{t+1}(c)=\underbrace{\frac{1}{2}(r_{t+1}-2c)}_{\text{fixed half exposure}}
+\underbrace{\delta_t r_{t+1}-2c\delta_t}_{\text{incremental gross return and cost}}.
\label{eq:returndecomp}
\end{equation}
This is an additive identity for daily returns and their arithmetic means, not for compounded wealth or Sharpe ratios. A negative incremental cost means reduced exposure relative to the benchmark; it does not mean a rebate. For two models A and B, the paired mean difference likewise satisfies
\begin{equation}
\Delta\bar R(c)=\Delta\bar G-2c\,\Delta\bar w.
\label{eq:paireddecomp}
\end{equation}
Equation~\eqref{eq:paireddecomp} permits a mean gross/net decomposition from mean weights and mean net differences alone. Applied to the aggregate inputs alone it cannot recover separate confidence intervals for the two components; the path-level records of Appendix~\ref{app:reconstruction} support them directly, and Section~\ref{sec:economics} reports them.

The 1 bps case is a low-friction sensitivity and 5 bps the prespecified higher-cost benchmark, not an estimate of institutional or retail execution costs. Auction costs depend on order size, timing, volatility and imbalance direction; a quoted spread alone is not an all-in execution cost.\footnote{NYSE, Choey Li (2023), ``Closing Auction: Immediate market impact, price drift and transaction cost of trading Part 2,'' 17 October. \url{https://www.nyse.com/data-insights-blog/closing-auction-immediate-market-impact-price-drift-and-transaction-cost-of-trading-part-2-eqjcqvixj8}. This corporate-stock evidence motivates cost heterogeneity; it does not calibrate ETF auction costs for this sample.}

Net returns are compounded. Sharpe uses $\sqrt{252}\,\bar R/s_R$ with zero risk-free return; drawdown includes initial capital. This conventional scaling is not serial-dependence adjusted \citep{Lo2002}. Cash and daily fully invested intraday exposure are timing-matched benchmarks.

\subsection{Nested benchmarks and two structural properties}\label{sec:benchmarks}
Let $N^{\mathrm{up}}_t$ and $N^{\mathrm{low}}_t$ count upper and lower destination labels among $s_{2\mid t},\ldots,s_{t\mid t}$, and let $\bar q_t=(N^{\mathrm{up}}_t+\alpha)/(N^{\mathrm{up}}_t+N^{\mathrm{low}}_t+2\alpha)$. The IID forecast of Equation~\eqref{eq:iid} has exactly this conditional upper share. Two benchmarks retain $\bar q_t$ and change only the middle probability:
\begin{equation}
p^{B}_t=\bigl((1-m^{B}_t)(1-\bar q_t),\;m^{B}_t,\;(1-m^{B}_t)\bar q_t\bigr),\qquad B\in\{\mathrm V,\mathrm M\}.
\label{eq:benchmarks}
\end{equation}
For V, $m^{\mathrm V}_t$ is the smoothed probability of a middle label after the current label in a two-state chain fitted to the indicators $\mathbf 1\{s_{u\mid t}=2\}$; it uses only whether sessions were middle or outer. For M, $m^{\mathrm M}_t=\widehat P_{s_{t\mid t},2,t}$, the unrestricted middle probability. Both use $\alpha=1$ and the same labelled history as the other models. By Equation~\eqref{eq:scorefactor}, IID, V and M share the conditional upper-versus-lower term, while M and the unrestricted forecast U share the middle-versus-outer term. For every target, therefore,
\begin{equation}
\ell^{\mathrm U}-\ell^{\mathrm{IID}}=(\ell^{\mathrm V}-\ell^{\mathrm{IID}})+(\ell^{\mathrm M}-\ell^{\mathrm V})+(\ell^{\mathrm U}-\ell^{\mathrm M}),
\label{eq:layers}
\end{equation}
where the first two layers involve only the middle-versus-outer term and the third only the conditional term. Because $R_{t+1}(c)$ is linear in the weight, the same telescoping holds exactly for net returns. The layers refine the factorization: $\ell^{\mathrm M}-\ell^{\mathrm{IID}}$ is the unrestricted middle-versus-outer difference and $\ell^{\mathrm U}-\ell^{\mathrm M}$ the conditional difference reported in Section~\ref{sec:components}. We call the three layers middle-versus-outer persistence (the persistence layer), additional middle information and conditional direction. The first records whether middle and outer tercile sessions tend to follow their own kind; it is related to volatility clustering but is not a volatility model. V and M were specified after the primary results were known.

\textbf{Property 1 (balanced marginal direction).} Suppose the cut points are the linearly interpolated empirical terciles of the same returns that are labelled and counted, with no duplicated return at a cut point. If $n$ returns are labelled, the number at or below the lower cut point exceeds the number above the upper cut point by zero or one, and dropping the first label changes the difference by at most one. Hence $|N^{\mathrm{up}}_t-N^{\mathrm{low}}_t|\le2$, so that $|2\bar q_t-1|\le2/(N^{\mathrm{up}}_t+N^{\mathrm{low}}_t+2\alpha)$ and, from Equation~\eqref{eq:weight},
\begin{equation}
\bigl|w^{B}_t-\tfrac12\bigr|=\tfrac12(1-m^{B}_t)\,|2\bar q_t-1|\le\frac{1}{N^{\mathrm{up}}_t+N^{\mathrm{low}}_t+2\alpha},\qquad B\in\{\mathrm{IID},\mathrm V,\mathrm M\}.
\label{eq:halfbound}
\end{equation}
The count statement follows from the interpolation rule: with $h=(n-1)/3$, exactly $\lfloor h\rfloor+1$ returns lie at or below the lower cut point and $n-\lfloor 2h\rfloor-1$ above the upper one. The primary and rolling rules satisfy the premise; fixed cut points do not. Any score improvement confined to the first two layers therefore leaves the allocation within this bound of one-half under the primary rule. Every final-period origin satisfies Equation~\eqref{eq:halfbound}; Appendix Table~\ref{tab:layerrules} reports the realized deviations.

\textbf{Property 2 (what the reversible projection changes).} For a positive transition matrix $P$ with stationary distribution $\pi$, let $F_{ij}=\pi_iP_{ij}$ and $S=(F+F^\top)/2$. The reversible kernel closest to $P$ in Kullback--Leibler divergence rate is the additive reversiblization $P^\ast_{ij}=S_{ij}/\pi_i$ \citep[Theorem~6.1]{WolferWatanabe2021}, which the reversible maximum-likelihood estimator targets as the sample grows. $P^\ast$ has the same stationary distribution. Because the flow out of the middle state equals the flow into it, $S_{21}+S_{23}=F_{21}+F_{23}$: the probability of leaving the middle state and the aggregate stationary flow in each direction between the middle state and the pooled outer states are unchanged. The projection changes two things only. It moves middle probability between the outer current states, $P^\ast_{12}-P_{12}=-J_{12}/(2\pi_1)$ and $P^\ast_{32}-P_{32}=J_{23}/(2\pi_3)$, with zero stationary-weighted net change, and it changes the upper-versus-lower split from every state. Detailed balance thus cannot alter the lumped middle-versus-outer flows that V estimates; any population effect on the score runs through the conditional direction and through the reallocation of middle probability between lower and upper current states. A finite-sample fit can still differ in both components.

The divergence rate $D=\sum_{ij}F_{ij}\log(F_{ij}/S_{ij})$ is the per-step expected log-loss cost of the projection. If both estimators are consistent for their targets, a first-order expansion gives
\begin{equation}
\mathrm E\bigl[\ell^{\mathrm R}-\ell^{\mathrm U}\bigr]\approx D-\frac{1}{2n},
\label{eq:refcost}
\end{equation}
where $n$ is the number of training transitions and $1/(2n)$ is the variance cost of the unrestricted model's one additional free parameter, the same term that underlies information-criterion order selection for finite Markov chains \citep{Tong1975}. The approximation assumes a stationary first-order chain with fixed labels, ignores smoothing, and treats the variance term as its value near reversibility. It implies an approximate crossover length $n^\ast=1/(2D)$ below which the restriction tends to help. It is a reference for magnitudes, not a test.

\subsection{Inference and its scope}
Paired non-circular moving blocks \citep{Kunsch1989} generate 5,000 resamples with seed 20261113. Block length is 20 sessions in the primary analysis and 10 or 60 in sensitivity checks. Within each comparison, the same sampled days are used for both models. The contrasts are unrestricted minus IID and reversible minus unrestricted; negative score differences favor the first model, whereas positive return and Sharpe differences favor it. Percentile intervals are pointwise unless explicitly labeled simultaneous. An interval containing zero does not establish equivalence.

Three families are distinguished: six log-loss contrasts across the three assets, six Brier contrasts, and an exploratory family of twelve economic contrasts comprising six mean net-return and six Sharpe differences. Each family is adjusted separately using aligned block draws across assets and contrasts (Section~\ref{sec:joint}; Appendix~\ref{app:families}). There is no adjustment across the three families. Component, calendar-segment and alternative-rule intervals remain exploratory and pointwise.

Resampling conditions on the saved forecast paths and does not refit thresholds or model parameters. It measures variation in the realized paired loss or return sequences under the dependence assumptions of block resampling, without regenerating estimation or design selection across new histories. Parameter estimation can affect forecast-evaluation inference \citep{West1996}. The present resampling does not assess the bias--variance trade-off across newly estimated forecast histories. It also does not establish that the reported intervals are necessarily narrower than unconditional intervals: estimation effects and their cross-covariances differ across models, so parameter count alone does not determine the direction of the change in uncertainty for a paired loss difference. No general coverage result is asserted for this expanding-window procedure with changing thresholds. Block-length sensitivity does not resolve structural change, and the related ETFs do not provide independent replications.

\section{Results and interpretation}
\subsection{Forecast accuracy}
Table~\ref{tab:scores} reports the lowest log loss and Brier score for the unrestricted chain on all assets. Reversible-minus-unrestricted log-loss differences are 0.001315, 0.000685 and 0.001595 for SPY, QQQ and IWM. All primary intervals contain zero (Table~\ref{tab:intervals}). Unrestricted-minus-IID intervals lie below zero pointwise for all three assets, while QQQ and IWM cross the 5\% threshold under the stepdown adjustment of Section~\ref{sec:joint}; IWM does not do so under the other reported familywise procedures. Related ETFs are not independent replications. Online Appendix A supplies development/validation context and full inference sensitivities. Relative to IID mean log loss, the unrestricted reductions are approximately 0.484\%, 1.065\% and 0.450\% (Table~\ref{tab:scale}). These are reductions in a scoring loss, not percentage-point gains in classification accuracy or investment returns. The intervals for reversible-minus-unrestricted differences neither demonstrate a reversible advantage nor exclude every practically relevant difference.

\subsection{Joint inference across the six log-loss contrasts}\label{sec:joint}
Table~\ref{tab:scores} and Table~\ref{tab:intervals} report six primary log-loss contrasts, two per asset. Read one at a time, three of them exclude zero. That reading ignores two things: the six statistics are computed on the same 2,766 sessions for three highly correlated funds, and separate 95\% intervals do not control the probability of at least one false exclusion.

We therefore test the six contrasts jointly. The daily paired loss differences are aligned on the common target sessions and resampled with the same moving-block indices in every draw, so that dependence across assets and across contrasts is preserved. Observed statistics and centered bootstrap draws use the same bootstrap standard deviation for each contrast (the standard error is not re-estimated within each draw), and two-sided adjusted $p$-values follow a Romano--Wolf stepdown procedure \citep{RomanoWolf2005}: the maximum statistic is first screened over all six hypotheses, and then over the hypotheses not yet rejected. A single-step max-$|t|$ critical value supplies simultaneous intervals; because it controls the family in one step it is more conservative than the stepdown test, and the two procedures need not agree at the boundary.

Table~\ref{tab:joint} gives the result. QQQ's unrestricted-minus-IID contrast is significant under all three reported procedures: stepdown-adjusted $p=0.008$, Holm-adjusted $p=0.010$, and a simultaneous interval of $[-0.0210,-0.0024]$. IWM is significant under stepdown ($p=0.048$) but not under Holm ($p=0.056$), and its single-step simultaneous interval contains zero. We therefore report a method-sensitive rejection for IWM, rather than treating all familywise procedures as yielding the same decision. SPY is not significant under either adjusted test ($p=0.093$ and $0.097$). None of the reversible-minus-unrestricted contrasts is significant; the smallest stepdown-adjusted $p$-value is 0.191. Familywise interpretation remains conditional on the validity of the underlying block-bootstrap approximation.

Two further families are reported in Appendix~\ref{app:families}: six Brier-score contrasts, and one exploratory economic family of twelve contrasts covering mean net returns and Sharpe ratios, each adjusted within itself. The Brier contrasts reject for QQQ (adjusted $p=0.006$) and IWM (0.041) but not SPY (0.067). In the economic family the Sharpe contrasts reject for SPY and IWM on both comparisons, while among mean net-return contrasts only IWM's unrestricted advantage survives (0.022). These are separate families, not additional corrections of the six log-loss contrasts, and the statements below about score advantage are statements about log loss.

The substantive ranking and the strength of evidence should be kept separate. Detailed balance has no demonstrated log-loss advantage under the reported procedures. The unrestricted model has the lowest point estimate for every asset; its advantage over IID has consistent familywise support for QQQ, method-sensitive support for IWM, and no rejection for SPY.

\subsection{Which score component carries the improvement}\label{sec:components}
Equation~\eqref{eq:scorefactor} splits each realized log loss into a middle-versus-outer term and a conditional upper-versus-lower term. Table~\ref{tab:components} reports the mean of each term in the saved final forecasts, its paired difference against IID, and a percentile interval from the same moving-block design used for the total loss.

The unrestricted advantage is concentrated in the first term for SPY and QQQ. Mean reductions are 0.00535 and 0.01163, and both pointwise intervals lie below zero: $[-0.00993,-0.00088]$ for SPY and $[-0.01850,-0.00456]$ for QQQ. Their conditional terms move by $+0.00004$ and $-0.00007$, both indistinguishable from zero. IWM is the exception. Its middle-versus-outer reduction of 0.00363 has an interval containing zero, while its conditional reduction of 0.00131 excludes zero and accounts for 26.5\% of its total reduction of 0.00494. Thus, the conditional component shows no detectable improvement for SPY or QQQ, while IWM has a pointwise improvement in that component.

The position rule makes that distinction concrete. With $w_t=\tfrac12+\tfrac12(1-m_t)(2q_t-1)$, the local sensitivities are $\partial w/\partial m=-(2q_t-1)/2$ and $\partial w/\partial q=1-m_t$. Across the unrestricted model's saved final forecasts, the mean of $|\partial w/\partial m|$ is approximately 0.015 and the mean of $\partial w/\partial q$ is 0.66, so the allocation is roughly forty times more sensitive to the conditional upper-versus-lower probability than to the middle-versus-outer one. These derivatives compare equal local changes in the two probability coordinates; their ratio is not an attribution of realized return gains. A middle-tercile score improvement can have little effect on exposure when $q_t$ is close to one-half; Section~\ref{sec:layers} replaces this local comparison with an exact one.

The reversible chain reproduces the first term within noise: differences from the unrestricted chain are $-0.00013$, $+0.00034$ and $-0.00008$, all intervals containing zero. It loses ground in the second term, by $+0.00144$, $+0.00035$ and $+0.00167$, with the SPY and IWM intervals above zero. The larger component differences occur in the conditional upper-versus-lower term, which also contributes to IWM's gain over IID. This alignment is descriptive: the conditional score is not a circulation statistic, and imposing detailed balance can change both components. The component intervals are pointwise and do not identify a causal mechanism.

\subsection{Benchmark layers: middle-versus-outer persistence and direction}\label{sec:layers}
Table~\ref{tab:layers} applies Equation~\eqref{eq:layers}. The persistence layer, V minus IID, reduces log loss by 4.36, 9.60 and $3.24\times10^{-3}$ nats for SPY, QQQ and IWM, with every interval excluding zero; this is 82\%, 82\% and 66\% of the unrestricted-minus-IID difference. The additional middle information in M adds 0.99, 2.02 and $0.39\times10^{-3}$, each interval containing zero. The direction layer is $+0.04$, $-0.07$ and $-1.31\times10^{-3}$ and excludes zero only for IWM. Relative to V, the unrestricted chain's remaining advantage is $-0.96$, $-2.10$ and $-1.70\times10^{-3}$, with intervals $[-3.75,+1.93]$, $[-4.56,+0.66]$ and $[-3.80,+0.32]$. A two-state chain that records only whether sessions were middle or outer therefore reproduces most of the score gain, and the remainder is not distinguishable from zero for any asset.

The allocation sees the layers very differently. By Property~1, the V and M allocations stay within $2\times10^{-4}$ of one-half at every origin, and their net-return layers are below $2\times10^{-5}$ basis points per day in absolute value. Essentially all of the incremental net return therefore arises in the direction layer: 0.0425, 0.0200 and 0.0728 basis points per day, with the intervals of Section~\ref{sec:economics}. The three assets illustrate the three possible combinations. For IWM the direction layer improves both the score and net return; for QQQ it improves neither detectably; for SPY it leaves the score unchanged ($+0.04\times10^{-3}$, interval $[-1.04,+1.05]$) while its pointwise net-return interval excludes zero. The 0.0425 basis-point SPY gain does not survive the twelve-contrast economic family ($p=0.184$).

Appendix Table~\ref{tab:layerrules} repeats the layers under the four rules of Section~\ref{sec:sensitivity}. The persistence share ranges from 56\% to 108\% across the twelve rule-by-asset combinations. The direction layer's log-loss interval excludes zero in one of twelve, the primary IWM case, whereas its net-return interval excludes zero in nine. Under fixed cut points Property~1 does not hold: the marginal upper share lies between 0.515 and 0.531 for SPY and QQQ, so every model carries a common tilt. Even then the first two layers contribute 15\%, 22\% and 2\% of the unrestricted-minus-IID net difference. These are pointwise intervals on overlapping designs, and the counts are descriptive rather than a test.

One explanation for the dissociation is that the allocation earns return magnitudes, whereas tercile scores record only the class. Weighting the conditional term on each outer-tercile day by $|r_{t+1}|$, normalized by its outer-day mean, makes all twelve direction estimates favour the unrestricted model, against nine of twelve unweighted, and increases the absolute ratio of estimate to bootstrap standard error in every case; three of twelve intervals exclude zero, all for IWM. This weighted quantity is a diagnostic, not a proper score for the tercile event, and the result does not separate the magnitude explanation from the low power of a binary component.

IID is also a weak reference for score levels. Realized class frequencies drift away from one-third because cut points come from the full history (Appendix~\ref{app:families}). A forecast that uses only the smoothed class frequencies of the latest 252 labelled sessions, with the primary cut points and no state dependence, has lower full-period point log loss than the unrestricted chain for all three assets, by 9.13, 10.18 and $2.43\times10^{-3}$ nats, with intervals containing zero (Appendix Table~\ref{tab:rollfreq}). This adaptation is unstable: with 1,260 sessions it beats the unrestricted chain in 2015--2018 for all three assets and loses to it in 2019--2022 for QQQ and IWM, each with an interval excluding zero. Their recent-window class counts are not balanced, so Property~1 does not apply to them, and their allocations are not evaluated. They show that time variation in scores is strongly affected by marginal drift that neither IID nor the Markov models track, and that the unrestricted chain's advantage over IID should not be read as an advantage over adaptive marginal forecasts. The window lengths were chosen after the primary results.

\subsection{What detailed balance removes: circulation flows}\label{sec:circulation}
For a stationary distribution $\pi$ and transition matrix $P$, the net probability flow from $i$ to $j$ is $J_{ij}=\pi_iP_{ij}-\pi_jP_{ji}$. Stationarity makes the three independent entries equal, $J_{12}=J_{23}=J_{31}=j_t$, so one scalar measures the mass rotating around the state cycle at each origin. Detailed balance is $j_t=0$, which the reversible parameterization imposes by construction. In this sample, the fitted stationary distribution stays within $4\times10^{-4}$ of uniform at every origin, and $j_t$ is then approximately one third of the corresponding pairwise difference in transition probabilities.

Figure~\ref{fig:circulation} plots $j_t$ for the three assets and Table~\ref{tab:circulation} summarizes it. The scalar is negative at every one of the 2,766 evaluation origins for every asset, with means of $-1.40$, $-0.77$ and $-1.32$ percentage points of daily flow. In the period-average matrices the three pairwise differences have nearly equal magnitude: for SPY, $P(\text{Up}\mid\text{Down})-P(\text{Down}\mid\text{Up})=+4.2$ and the two other pairs are $-4.2$ percentage points; the corresponding magnitudes are 2.3 for QQQ and 4.0 for IWM. The reversible fit imposes $\pi_iP_{ij}=\pi_jP_{ji}$. With the estimated stationary distribution that close to uniform, its fitted entries are within $3\times10^{-4}$ of pairwise averages, and it shifts individual off-diagonal entries by 1.2 to 2.1 percentage points. That near-equality of entries is a property of this estimate, not of the estimator: when the stationary distribution is not uniform, detailed balance does not imply $P_{ij}=P_{ji}$.

The series is highly persistent, with first-order autocorrelation between 0.990 and 0.997. The agreement in sign across origins is therefore not 2,766 independent confirmations; it describes one persistent feature of this sample, and the annual means move by at most a few tenths of a point. The reversible-minus-unrestricted point estimates range from 0.0007 to 0.0016 nats of log loss; each asset-specific confidence interval contains zero (Table~\ref{tab:intervals}). The fitted chain carries the asymmetry at every origin and the restriction removes it by construction, but whether the population chain does so cannot be settled here: an expanding-window path reuses almost all of its history between adjacent origins, so persistence is expected even if the underlying circulation is zero. No test of $j=0$ in the population is performed. Section~\ref{sec:simulation} quantifies the cost of the restriction under generating processes whose circulation is known; that explains a mechanism under stated assumptions and does not identify it in these data.

Equation~\eqref{eq:refcost} gives a first-order reference for the size of the observed contrasts (Appendix Table~\ref{tab:restriction}). Averaging $D_t-1/(2n_t)$ over the fitted matrices gives 1.289, 0.355 and $1.131\times10^{-3}$ nats for SPY, QQQ and IWM, against observed reversible-minus-unrestricted differences of 1.315, 0.685 and $1.595\times10^{-3}$. These matrices incorporate earlier final-period outcomes. Using only the matrix of the first final-period origin, estimated from data through 31 December 2014, gives 1.087, 0.564 and $0.650\times10^{-3}$. Every reference lies inside the corresponding interval and amounts to between 0.7 and 1.4 bootstrap standard errors of the contrast. Under this approximation, an interval containing zero is the more likely outcome in a sample of this length even if the fitted circulation were the population value; the non-rejection is uninformative about whether detailed balance holds. The references assume a stationary first-order chain with fixed labels and were formulated after the results; they are not independent predictions.

Consistent with Property~2, the reversible fit departs from the unrestricted fit mainly in the conditional term for SPY and IWM: the middle-versus-outer differences are $-0.13$ and $-0.08\times10^{-3}$ nats, against $+1.44$ and $+1.67\times10^{-3}$ in the conditional term. For QQQ the two departures are similar, $+0.34$ and $+0.35\times10^{-3}$, consistent with the projection's reallocation of middle probability between outer current states. As a further implementation check, the closed-form symmetric estimator $(C_{ij}+C_{ji})/(r_i+c_i)$, with smoothed counts $C$ and their row and column sums $r$ and $c$, maximizes the reversible likelihood whenever row and column sums are equal; here they differ by at most one transition, and the closed form agrees with every saved reversible forecast to within $1.4\times10^{-4}$.

\subsection{Absolute performance and incremental economic value}\label{sec:economics}
All nine allocations lose capital at primary costs (Table~\ref{tab:economics}; Figure~\ref{fig:wealth}). SPY unrestricted cumulative return falls from 31.21\% gross to $-0.47\%$ at 1 bps and $-67.05\%$ at 5 bps per side. QQQ yields 41.66\% gross and 7.40\% at 1 bps, but a net loss of 64.52\% at 5 bps per side; IWM already loses 18.03\% gross (Table~\ref{tab:costs}). These are full-period cumulative returns, not annualized rates.

Across saved final forecasts, all allocations lie between 0.4708 and 0.5273 (Online Appendix Table A5). This is a descriptive audit of frozen predictions. The rule therefore implements a small probability-driven adjustment around a persistent half exposure. Daily liquidation and re-entry produce $T_{t+1}\approx1$, charging approximately 5 bps of portfolio capital per day at the primary cost.

The scale of this drag can be seen without a market-data model. If exposure were exactly one-half and every underlying intraday return were zero, wealth after 2,766 sessions at $c=0.0005$ would be $(1-c)^{2766}\approx0.250738$, a loss of 74.93\%. This is a counterfactual accounting illustration, not an attribution of 74.93 percentage points of the observed losses. Actual wealth also reflects return realizations, changing exposure and compounding. It explains why absolute losses under the fixed rule must be separated from the incremental value of its probability adjustment.

Table~\ref{tab:scale} applies Equation~\eqref{eq:paireddecomp} to unrestricted-minus-IID mean net differences and retained mean weights. The approximate daily net gains are 0.0425 bps for SPY, 0.0200 bps for QQQ and 0.0728 bps for IWM. Their implied gross gains are 0.0381, 0.0249 and 0.0703 bps; mean cost changes are $-0.0045$, 0.0049 and $-0.0025$ bps. Simple annualization of the mean net differences gives 0.1071, 0.0504 and 0.1834 percentage points per 252 sessions. These small relative improvements coexist with negative absolute performance. They are arithmetic rescalings of retained point estimates, not new hypothesis tests or compound annual returns.

Applying the same moving-block design to the component series, rather than to the aggregates, gives separate intervals. The incremental gross return averages 0.038, 0.025 and 0.070 basis points per day, with intervals $[-0.003,0.083]$, $[-0.021,0.073]$ and $[0.025,0.121]$. The net difference averages 0.043, 0.020 and 0.073 basis points with intervals $[0.001,0.087]$, $[-0.026,0.069]$ and $[0.028,0.123]$. Only the IWM gross interval and the SPY and IWM net intervals exclude zero. The same resampling gives the cost term itself: $-0.0045$, $+0.0049$ and $-0.0025$ basis points per day, with intervals $[-0.0090,0.0002]$, $[-0.0009,0.0107]$ and $[-0.0084,0.0030]$. The absolute mean cost difference is below 0.005 basis points for each asset. At these point estimates, the gross difference is positive and larger in magnitude than the cost difference; this accounting comparison is not a causal attribution. Interval endpoints are not subtracted from one another to obtain a cost interval, which would be invalid. For QQQ both the gross and the net differences are indistinguishable from zero at this precision, consistent with Appendix Table~A3. These component intervals are pointwise; in the separate twelve-contrast economic family, only IWM's unrestricted-minus-IID mean net difference remains significant.

Figure~\ref{fig:calibration} shows bin-average probabilities, and Section~\ref{sec:components} attributes the score gain to a component rather than to the total. The allocation ignores within-state return magnitudes, so aggregate state scores alone cannot determine an economic break-even cost. Such a break-even would divide incremental gross return by incremental turnover; the two models differ by at most $1.0\times10^{-3}$ round trips per session in mean turnover, the implied ratio changes sign across assets, and it is therefore not reported. At 1 bps per side, unrestricted SPY and IWM return $-0.47\%$ and $-37.83\%$, whereas QQQ returns 7.40\% (Table~\ref{tab:costs}); profitability remains cost- and asset-dependent. No alternative allocation is selected from these outcomes.

IID exposure is nearly constant, ranging from 0.499864 to 0.500000 across the three assets. For an exactly constant $w>0$, $R_{t+1}=w(r_{t+1}-2c)$ implies $\operatorname{Sharpe}(wX)=\operatorname{Sharpe}(X)$, since the mean and standard deviation scale together. This explains IID and fully invested intraday Sharpe ratios agreeing to three decimals in Table~\ref{tab:economics}; the observed near-constant weights do not imply exact equality beyond rounding, equal drawdowns or equal cumulative returns. On the reconstructed panel the IID mean net return differs from an exact fixed-half benchmark by at most $6\times10^{-8}$ per session, so the IID comparison is also, to that precision, a comparison against fixed half exposure.

QQQ provides a particularly clear contrast. Its unrestricted-minus-IID log-loss gain is largest ($-0.011697$, versus $-0.005314$ for SPY and $-0.004942$ for IWM), yet its primary mean-net-return and Sharpe-difference intervals contain zero; the latter is $[-0.008712,0.017061]$ (Online Appendix Table A3). SPY and IWM have positive primary relative-Sharpe intervals (Tables A2 and A4), despite smaller score gains. These conditional, pointwise comparisons illustrate that score improvements need not rank economic improvements monotonically; they are not a formal test of cross-asset differences, and positive relative Sharpe does not imply positive absolute profitability. Section~\ref{sec:layers} locates the reason: QQQ's larger score gain is mostly middle-versus-outer persistence, which cannot move the allocation.

Appendix Table~\ref{tab:statemeans} summarizes the same final-period data by the state of the origin session. Mean next-session returns after a lower-tercile session are 7.61, 13.37 and 7.64 basis points. For QQQ the pointwise interval is $[5.73,22.53]$ and the nine-cell simultaneous interval $[1.68,25.05]$; both contain the 10 basis-point round trip at the primary cost. The unrestricted allocation adds 1.67, 0.28 and 1.50 percentage points of exposure after lower-tercile sessions, and QQQ's highest mean exposure follows middle sessions (51.04\%). Probability tilts and conditional mean returns are different objects: terciles discard magnitudes, and the rule maps class probabilities rather than expected returns into exposure. These are in-sample conditional averages. Choosing participation states from them and reporting returns in the same period would be a selected rule, and no such rule is evaluated here.

For SPY, the unrestricted-minus-IID Sharpe difference is significant in the twelve-contrast economic family ($p=0.032$), whereas the mean net-return difference is not ($p=0.184$), but the unrestricted daily net-return standard deviation is slightly higher: 0.00425030 versus 0.00423084, an increase of 0.46\%. With both mean returns negative, this larger denominator makes the conventional Sharpe ratio less negative; an exact symmetric decomposition assigns 0.015912 of the 0.022689 point-estimate improvement to the mean change and 0.006777 to the inverse-standard-deviation change.\footnote{For $S_M=\sqrt{252}\mu_M/\sigma_M$ and $M\in\{U,I\}$, the symmetric identity is $S_U-S_I=\sqrt{252}\{(\mu_U-\mu_I)(\sigma_U^{-1}+\sigma_I^{-1})/2+(\mu_U+\mu_I)(\sigma_U^{-1}-\sigma_I^{-1})/2\}$. Here $I$ denotes IID and $\sigma_M$ is the sample daily net-return standard deviation. The two terms are descriptive accounting contributions, with no separate significance tests.} These contributions describe the change in the ratio, not the cause of its significance: the paired ratio statistic has a different bootstrap distribution from the mean contrast, and neither the different $p$-values nor the decomposition demonstrates risk reduction or volatility timing.

\subsection{Threshold and window rules}\label{sec:sensitivity}
The primary design re-estimates terciles from every past session at each origin. Table~\ref{tab:sensitivity} repeats the final-period comparison under four alternatives: cut points fixed at their first eligible values, rolling windows of 2,520 and 1,260 sessions applied to both the cut points and the state history, and a common start date of 26 May 2000 for all three assets.

Two features survive every rule. In all fifteen asset-by-rule combinations the average unrestricted-minus-IID loss difference is negative, between $-0.0017$ and $-0.0148$ nats, and the reversible-minus-unrestricted difference is positive. These are directions in the listed designs, not dominance relations between models. For reversible minus unrestricted, the only interval that excludes zero is IWM under fixed cut points, and it favors the unrestricted chain. The rows share sessions and assets; IWM's common-start row reproduces its primary row by construction.

Magnitude and precision depend on the rule. The unrestricted-minus-IID interval excludes zero under the primary design and under fixed cut points for all three assets, but under a 2,520-session rolling window it includes zero for QQQ, and under a 1,260-session window it includes zero for all three assets. The 1,260-session point estimates are 67\%, 60\% and 51\% smaller in magnitude than the primary ones, so the loss of precision reflects a smaller effect as well as a shorter estimation history. Shortening SPY's history to a common 2000 start changes its point estimate from $-0.00531$ to $-0.00536$ and widens the interval to include zero. The negative full-period point estimates are robust in direction across these rules, but their precision depends on the estimation history. This design sensitivity is separate from stability across calendar periods.

Absolute log-loss levels refer to different target definitions across rules. With fixed cut points, unrestricted loss falls from 1.0933 to 1.0732 for SPY and from 1.0870 to 1.0306 for QQQ. These reductions can partly reflect a more imbalanced class distribution; they cannot be interpreted as improved forecasting of a common target. Paired differences within each rule remain interpretable as model comparisons for that target definition. Their signs and sensitivity can be compared across rules, but their magnitudes do not measure gains on an identical prediction task.

The economic comparison behaves the same way. The mean net difference stays positive under every rule, between 0.02 and 0.14 basis points per day, and its interval contains zero in four of the fifteen rows while excluding it in the other eleven. These are pointwise comparisons across overlapping designs, rather than a family-level result.

Under fixed and rolling rules, the original optimizer reaches a line-search precision limit at 0.1\%--4.6\% of estimation origins. The v13 reproduction handles these cases by optimizing the same smoothed likelihood over simplex flow masses and polishing the stationarity equations. Acceptance requires positive finite flows and a projected-gradient infinity norm at most $10^{-8}$; this replaces the earlier softmax-gradient acceptance rule. The primary rule has no fallbacks. The saved v13 results reproduce the reported sensitivity estimates to their displayed precision; Appendix B distinguishes this numerical reproduction from statistical validation.

\subsection{Time stability within the final period}\label{sec:stability}
Table~\ref{tab:segments} divides the final period into three calendar segments selected for this retrospective revision, and reports both contrasts within each. The unrestricted advantage is not uniform. It is largest in 2015--2018, when all three intervals exclude zero ($-0.0118$, $-0.0299$ and $-0.0079$ for SPY, QQQ and IWM), weakens in 2019--2022, and in 2023--2025 changes sign for SPY and QQQ ($+0.0046$ and $+0.0111$, the latter with an interval excluding zero) while remaining essentially zero for IWM. The QQQ gap in the first segment is two and a half times its full-period estimate. The persistence layer of Section~\ref{sec:layers} carries most of this variation: it accounts for 70\%, 77\% and 59\% of the 2015--2018 gains, and for QQQ in 2023--2025 it contributes $+8.88\times10^{-3}$ nats, with interval $[+0.36,+17.83]$, to the $+11.06\times10^{-3}$ reversal. The middle-versus-outer persistence learned from the expanding history thus helped early in the evaluation and hurt QQQ late in it.

Figure~\ref{fig:rolling} shows the same information as a 252-session rolling mean. The paths cross zero around 2022--2023 for all three assets and stay near or above it afterwards. Overlapping windows are not independent observations and the path is descriptive; what it shows is that the full-period average is a weighted compromise between a favourable early period and an unfavourable recent one rather than a stable property of the model.

The reversible-minus-unrestricted difference is positive in all nine segments, with one interval (IWM, 2015--2018) excluding zero and lying on the worse side for the reversible model. The point-estimate sign of the restriction contrast is more consistent across segments than that of unrestricted minus IID; the segment-specific intervals do not establish a stable population disadvantage or test equality across segments.

\subsection{Controlled simulation of the restriction}\label{sec:simulation}
Table~\ref{tab:simulation} examines a designed question: when observations follow a stationary three-state chain with known circulation, how does imposing detailed balance affect expected log loss at different training lengths?

The baseline transition matrix $P_0$ has diagonal entries 0.35, off-diagonal entries 0.325 and stationary distribution $\pi_i=1/3$. Its stationary flow matrix is $S=P_0/3$. The generating flow is $F=S+(j/2)C$, where $C_{12}=C_{23}=C_{31}=1$, the reverse entries are $-1$, and the diagonal is zero. Setting $P=3F$ gives circulation $J_{12}=J_{23}=J_{31}=j$. The scenarios use $j=0,-0.007,-0.014$, equivalent to 0, $-0.7$ and $-1.4$ percentage points of daily flow. For each training length from 250 to 5,000 states, 500 chains are simulated and both estimators are fitted with the same smoothing. Expected log loss is evaluated directly against the generating matrix, with stationary row weights; no finite test path is drawn.

When detailed balance is correct, the restricted estimator has lower expected loss at every listed length. Its advantage decreases from $2.09\times10^{-3}$ nats at 250 states to $1.01\times10^{-4}$ at 5,000. Under misspecification, the relative benefit of variance reduction diminishes with training length and the restriction bias becomes dominant. For $j=-0.014$, the reversible-minus-unrestricted difference rises from $-0.638\times10^{-3}$ at 250 states to $+0.950\times10^{-3}$ at 1,260, $+1.165\times10^{-3}$ at 2,520 and $+1.257\times10^{-3}$ at 5,000. For $j=-0.007$, the sign changes between 1,260 and 2,520 states.

The first-order reference of Equation~\eqref{eq:refcost} reproduces all fifteen cells within 1.27 Monte Carlo standard errors (Appendix Table~\ref{tab:restriction}). For the two misspecified scenarios it implies crossover lengths $n^\ast$ of about 1,474 and 368 transitions, matching the sign changes in the table. The simulation therefore confirms a first-order approximation numerically; it adds no independent evidence about the data.

At 5,000 states, the stronger-circulation scenario produces a cost of approximately $1.3\times10^{-3}$ nats, of the same order as the empirical contrasts. This is an order-of-magnitude comparison: the simulation grid does not include the exact initial empirical training lengths of 5,522, 3,980 and 3,672 observations, and the weaker scenario is closer to QQQ's fitted circulation. The scenarios were selected after seeing the empirical results and omit the changing-threshold layer. They illustrate the bias--variance trade-off under stated generating processes; they neither identify the source of the observed ranking nor independently confirm it.

\subsection{Interpretation of the restriction}
The score decomposition and benchmark layers identify which forecasting errors change and whether the allocation can respond; the projection properties show which of those changes detailed balance can affect; the circulation path and the closed-form reference describe the fitted flow the restriction removes and the size of loss its removal implies; and the simulation checks that reference under known truth. Together they place the restriction's effect in the conditional direction and in the reallocation of middle probability between outer states, the same layer that carries the incremental economic return, and they show that its expected size is about one standard error of the observed contrast. None of these analyses establishes that population irreversibility causes the empirical score gap. The highly persistent fitted circulation reuses an expanding history, the primary reversible-minus-unrestricted log-loss intervals include zero, and the simulation was calibrated retrospectively. Likewise, the overnight-return literature does not explain IWM's intraday losses without an overnight decomposition or a causal design. The computational checks in Appendix B assess implementation and numerical agreement, not empirical power or economic identification.

\section{Limitations and conclusion}
The evidence is conditional on three related US equity ETFs, three return states, first-order dependence and the stated estimation rules. Changing terciles discard return magnitudes and relabel history across origins. Alternative thresholds and windows preserve the direction of the full-period comparisons but reduce their magnitude and precision in several cases. Calendar segments show a more substantive limitation: the unrestricted advantage is concentrated in the earlier period and reverses for SPY and QQQ in 2023--2025. The full-period mean should not be interpreted as a stable forecasting edge.

Data and execution assumptions also limit interpretation. Retrospective vendor adjustments are not a point-in-time feed. A common same-session adjustment factor cancels in $C/O$, but revisions to that ratio can change labels and fitted transitions. The cost scenarios are not estimates of actual ETF execution costs; fills, cash interest and exit notional are simplified. The daily-return decomposition is exact under this accounting but is not an additive decomposition of compounded wealth or Sharpe ratios.

Inference is conditional on saved forecasts, with unproven coverage under structural change and conventional Sharpe scaling that does not correct for serial dependence. Separate adjustments cover the log-loss, Brier and economic families; they do not control multiplicity across all analyses. Component, segment and alternative-rule intervals remain pointwise. The calibration simulation imposes independent draws from a frozen forecast sequence and therefore is not a general test of sequential calibration. The internal design freeze is not external preregistration, and retrospective diagnostics do not supply new out-of-sample evidence. The V and M benchmarks, rolling-frequency forecasts, magnitude-weighted direction term, segment boundaries and closed-form references were all specified after the primary results; the twelve rule-by-asset counts are descriptive. Property~1 depends on re-estimating terciles from the labelled history and fails under fixed cut points. The closed-form reference is a first-order approximation for a stationary chain with fixed labels. Any selection of a new allocation or estimation rule requires a separately declared evaluation.

Within these limits, the unrestricted model achieves modest average score gains over IID. Log-loss evidence is consistent across the reported familywise procedures for QQQ, sensitive to the procedure for IWM, and not significant for SPY after adjustment. Between 66\% and 82\% of each gain is middle-versus-outer persistence that a two-state chain reproduces, and the remainder is not distinguishable from zero relative to that chain. The advantage is strongest in 2015--2018 and reverses for SPY and QQQ in 2023--2025, for QQQ mainly through the persistence layer. All primary reversible-minus-unrestricted log-loss intervals contain zero; the exploratory fixed-cutpoint and early-period IWM exclusions favor the unrestricted model. Their point estimates are of the size a first-order projection cost implies, so the data neither demonstrate a forecast advantage from detailed balance nor have power to detect a cost of that size reliably; non-rejection does not establish equivalence.

The economic conclusion is specific to the decision rule. Daily liquidation makes maintaining near-half exposure costly even with uninformative probabilities. Total costs of about five basis points per day coexist with unrestricted-minus-IID net improvements of 0.02--0.07 basis points per day, but this scale comparison must not be read as the incremental cost of choosing the unrestricted forecast: mean cost differences between the two models are below 0.005 basis points per day. Absolute losses reflect realized gross returns, the baseline cost and compounding. The incremental return over IID arises almost entirely in the conditional direction layer, while the larger score layer moves the allocation only slightly under the primary rule. Absolute losses do not establish that a turnover-aware alternative would fail.

The contribution is an exact accounting for modest forecasting gains: nested benchmarks that split a score difference and a net-return difference into the same layers, properties of the restriction and of the labelling that say which layers each can affect, a closed-form reference for the size of the restriction's cost, and the separation of baseline exposure from incremental return and cost. In this application it shows that most of the score gain is a reversible, allocation-neutral layer, and that the economically relevant layer is the one in which scores have the least power. Transfer to another score, rule, state definition or holding period requires its own mathematical and empirical assessment.

\section*{Data and Code Availability Statement}
The replication supplement contains analysis code, aggregate results and SHA-256 fingerprints; third-party prices and row-level price-derived outputs are not distributed. Two forms of reproducibility are distinguished. Internal bit-exact auditability concerns rerunning the specified code on the same locally retained reconstructed price inputs: the v13 record reports identical forecast-file hashes and ten identical generated table files, which were checked against the existing local run artifacts for this revision. These checks do not make the inputs publicly available or establish independent replication.

External functional reproducibility concerns approximate agreement after a new authorized retrieval of the vendor series. The supplied reconstruction matches the archived row counts and dates, with published log loss reproduced to six decimals for SPY and IWM and within $3\times10^{-6}$ for QQQ; the QQQ unrestricted-minus-IID contrast differs by $2\times10^{-6}$ (Appendix~\ref{app:reconstruction}). A future download is not guaranteed to reproduce those tolerances or hashes. Common same-session adjustments cancel in $C/O$, whereas revisions to that ratio can change thresholds and labels. The archived snapshot is not distributed, and the replication code is available at \url{https://github.com/qqj53333/detailed-balance-etf}. The supplement records which checks were rerun and which verification claims are inherited from earlier records. The v16 benchmark layers, closed-form references and their tables are computed by \texttt{layers.py}, \texttt{restriction\_cost.py} and \texttt{tables\_v16.py} from the saved v13 forecasts and the same local price inputs, without refitting the reported models; they reproduce every saved unrestricted forecast to within $10^{-16}$ and every saved target label exactly. In v17 a separate implementation rebuilt all 8,298 final-period forecasts from the retained reconstruction, using reversible likelihood score equations, and checked next-session dates, labels, losses and net returns. The largest reversible probability discrepancy was below $1.3\times10^{-10}$. The public supplement includes this check and commands for replay with authorized price inputs.

\ifblind\else
\section*{Funding}
This research received no external funding.
\section*{Declaration of competing interest}
The author declares no financial or non-financial conflicts of interest.
\fi
\section*{Declaration of Generative AI in Scientific Writing}
Generative AI systems were used throughout this study for data work, code, statistical analysis and manuscript text, not only for language editing. Their roles are listed by system. The author directed the work, reviewed the outputs and takes full responsibility for the content.

For v17, Codex wrote and ran the separate implementation, checked the primary contrasts, and edited the manuscript and replication instructions. OpenAI ChatGPT and Codex assisted with literature searching, the implementation and tests of the frozen code package, and early manuscript drafts. Separate ChatGPT sessions prepared revision 12 and, after a peer-style review by Google Gemini, revision 14.

DeepSeek Harness, a DeepSeek-based coding agent, reconstructed the daily price panel from a public vendor endpoint, reran the unmodified frozen code, and wrote the analyses reported in Sections~\ref{sec:joint}, \ref{sec:components}, \ref{sec:circulation}, \ref{sec:sensitivity}, \ref{sec:stability} and \ref{sec:simulation} and Appendices~\ref{app:families} and~\ref{app:reconstruction}: the score decomposition and allocation sensitivities, circulation paths, the independent estimator check, the alternative estimation rules, familywise adjustments for three contrast families, calendar-segment and rolling stability, the controlled simulation and the calibration summaries. It also wrote the reproduction package, the generators for the corresponding tables, and the manuscript text of revisions 7--11, 13 and 15.

Anthropic Claude wrote the analyses added in revision 16, namely Sections~\ref{sec:benchmarks} and~\ref{sec:layers}, the closed-form references in Sections~\ref{sec:circulation} and~\ref{sec:simulation}, the conditional returns in Section~\ref{sec:economics} and Appendix~\ref{app:layers}, together with the scripts and table generators that compute them and the corresponding manuscript text.

Later AI-assisted revisions found and corrected errors introduced by earlier AI-assisted text, including a reversed sign in one score-component contrast, a calibration point estimate taken from bootstrap draws instead of the observed sample, and solver documentation that no longer matched the code. Verification is computational and is recorded in Online Appendix B and the replication supplement: the forecast panels reproduce with identical SHA-256 hashes from the local price files, the generated tables reproduce byte for byte, a 231-item automated comparison returns no mismatches, and revision-specific check scripts link reported numbers to computed outputs. These checks were written and run by the AI systems; the author reviewed their records. Figures are deterministic plots of computed outputs; no generative image synthesis or alteration of observations was used. The author chose the research question and design, decided which analyses and interpretations to retain, and is responsible for the accuracy of every claim.
